\subsection{Positive duality and the radical at the collision}
\begin{finding}[\fid{F47}: a positive adjoint realization must lose the radical]\label{f:positivecollision}
The actual algebra
\[
 A_1=\C\times\C\times M_2(\C[\varepsilon]/\varepsilon^2)
\]
cannot be the complexified Grothendieck ring of a multifusion
category. The same conclusion applies to a graded categorification
whose specialization at $v=1$ is explicitly required to have that form.
It is not asserted that forgetting any grading automatically produces
such a category.

More directly, if $\rho:A_1\to\End(H)$ is a representation on a
finite-dimensional positive-definite complex Hermitian space and
both named generators $\rho(P_1),\rho(P_2)$ are self-adjoint, then
\[
 0\times0\times M_2(\varepsilon\C)\subseteq\ker\rho.
\]
Thus this positive adjoint realization factors through the actual
rank-six classical quotient and cannot retain the full rank-ten
collision algebra faithfully.

On the regular rank-two center module, the natural trace pairing
is positive away from the collision and has determinant $\Delta^2$,
of order four at $v=1$. This is the smallest possible determinant
order of a regular symmetric form making the center coordinate
self-adjoint and positive at nearby real noncollision parameters.
The actual rank-64 tensor module, which kills the radical after
specialization, does have a positive form regular through the collision.
\end{finding}

\begin{proof}
The categorical input is the known finite-based-ring theorem:
the Grothendieck ring of a multifusion category is based, and its
complexification is semisimple
\cite[Definition~1.42.2 and Proposition~1.42.4, pp.~76--77;
Corollary~1.43.5, p.~80]{EGNO2009}.
The nonzero square-zero ideal
$J=0\times0\times M_2(\varepsilon\C)$ is the Jacobson radical
of $A_1$, by \fid{F19}. This contradicts the proposed semisimplicity
and proves the first assertion. The known theorem concerns finite
based rings, not every ring with nonnegative structure constants.

For the operator assertion put $e=P_1/4$, $f=P_2/4$ and
$z=(e-f)^2$. Finding~\fid{F19} shows that $z$ is central, zero
on the two scalar factors, and $3/4+\eta$ on the matrix factor,
where $\eta$ is a nonzero central square-zero element generating
$J$ as an ideal. The real polynomial
\[
 p(T)=1-\frac{(T-3/4)^2}{(3/4)^2}
\]
makes $p(z)$ the central projector onto the matrix factor.
Consequently $n=p(z)(z-3/4)$ is zero on the scalar factors,
equals $\eta$ on the matrix factor, satisfies $n^2=0$, and
generates $J$. Since $\rho(e)$ and $\rho(f)$ are self-adjoint,
$\rho(z)$ and then $\rho(n)$ are self-adjoint. For every $u\in H$,
\[
 \|\rho(n)u\|^2=\langle u,\rho(n)^2u\rangle=0.
\]
Thus $\rho(n)=0$ and $J\subseteq\ker\rho$. The quotient is
$\C\times\C\times M_2(\C)\simeq\C[S_3]$, with the actual
specialized generator map already identified in \fid{F19}.

For the pairing assertion work over
$R=\R[v,v^{-1}]_{(v-1)}$, set $t=v-1$, and write
\[
 C=R[x]/(x(x-\Delta)),\quad x=z_1-z,\qquad
 \Delta=\frac{D-1}{D^2}(v-v^{-1})^2=t^2u(t),\quad u(0)=3/4.
\]
The regular algebra trace gives, in the basis $(1,x)$,
\[
 G_{\rm tr}=\begin{pmatrix}2&\Delta\\\Delta&\Delta^2\end{pmatrix},
 \qquad \det G_{\rm tr}=\Delta^2.
\]
Indeed the traces of multiplication by $1,x,x^2$ are
$2,\Delta,\Delta^2$. The leading principal minor is 2 and the
determinant is positive at each nearby real point with
$\Delta\ne0$, proving positive definiteness and order four.

To prove optimality, right multiplication by $x$ has row matrix
\[
 X=\begin{pmatrix}0&1\\0&\Delta\end{pmatrix}.
\]
For a symmetric regular form
$G=\left(\begin{smallmatrix}a&b\\b&c\end{smallmatrix}\right)$,
the relation $XG=GX^{\mathsf T}$ is precisely $c=\Delta b$.
Positive definiteness for small real $t>0$ gives
\[
 a(t)>0,\qquad 0<b(t)<a(t)\Delta(t),\qquad
 \det G=b(a\Delta-b)>0.
\]
Let $s=\operatorname{ord}_t a$ and $r=\operatorname{ord}_t b$;
both are finite nonnegative integers, since generic positivity
excludes their identically vanishing. The inequality forces
$r\ge s+2$: otherwise $b/(a\Delta)$ tends to positive infinity
as $t\to0^+$, contradicting its being less than one. Therefore
\[
 \operatorname{ord}_t\det G
 \ge r+(s+2)\ge2s+4\ge4.
\]
The trace form attains the lower bound. The identical proof for a
positive $\Delta=t^m u(t)$ on a punctured real half-neighborhood
gives the optimal bound $2m$.

Finally, on either actual right block $N(a)$ of \fid{F21}, where
$a=\delta$ or $a=D-2$, put
\[
 G_a=\begin{pmatrix}1&a/D\\a/D&1\end{pmatrix}.
\]
Direct multiplication with
$P_1=\left(\begin{smallmatrix}D&0\\a&0\end{smallmatrix}\right)$
and $P_2=\left(\begin{smallmatrix}0&a\\0&D\end{smallmatrix}\right)$
gives $P_iG_a=G_aP_i^{\mathsf T}$. For every real $v>0$,
$\delta\ge2$, $D\ge4$, and
\[
 0<\delta/D=1/\delta<1,\qquad 0<(D-2)/D<1.
\]
Thus both forms are positive and specialize to
$\left(\begin{smallmatrix}1&1/2\\1/2&1\end{smallmatrix}\right)$.
Give the scalar sectors positive standard forms and take the
direct sum with multiplicities $20,4,16,4$. This is the asserted
positive form on the actual tensor module. The change of basis
of \fid{F21} is regular and invertible at $v=1$, so the form
transports to its original tensor basis without a pole there.
Its special action kills $J$, exactly as the operator theorem requires.
\end{proof}

An indefinite form can have nonzero self-adjoint nilpotents,
and purity alone imposes no positive-adjoint condition on the
named generators. A smooth proper dg category may also have
this algebra as the endomorphism algebra of one object, as
\fid{F40} and \fid{F44} show: this is a different role from the
whole based Grothendieck algebra. The result excludes the
specified finite rigid and faithful positive-adjoint routes,
not those other realizations or VIII's weaker signed-factor
positivity convention.

% Separate paragraph for Part III.
\paragraph{\fid{F47} in Grothendieck's language.}
The finite-flat scheme $\Spec C\to\Spec R$ has an algebra trace,
and $C\to C^\vee$, $a\mapsto(b\mapsto\Tr(ab))$, is a morphism
of finite locally free sheaves. Its determinant ideal is
$(\Delta^2)$ and its special-fiber kernel is the nilpotent
direction, by the displayed complete computation. Equivariant
modules use SGA~3, I, 4.7.2--3 and cote~90, scan~17.
The $K_0$ assertion applies the precisely cited modern theorem
on finite based rings; it is stronger than the virtual
$\lambda$ identities in SGA~6, V, \S\S2--3 and cote~133,
scan~18. For a proposed motivic polarization, the exact
additional assumptions are a positive-definite realization and
self-adjoint correspondences for the two named generators.
Under these assumptions the radical is necessarily killed.
Cote~16, scans~35,37,39 and~48--49, anchors the primitive-projector,
and conditional semisimplicity vocabulary. No
existence of such GCT correspondences or polarization is
attributed to those leaves. Neither a cohomological sign nor
Weil purity alone supplies this positive-adjoint structure.

% New bibliography keys:
% EGNO2009: P. Etingof, S. Gelaki, D. Nikshych, V. Ostrik,
% Tensor Categories, MIT 18.769 lecture notes (2009),
% Definition 1.42.2, Proposition 1.42.4, pp.76--77;
% Corollary 1.43.5, p.80; https://math.mit.edu/~etingof/tenscat1.pdf.
